\chapter{Calcul de prix et méthodes numériques}\label{ch:iv2:pricing-and-numerics}

\begin{ivtrains}
Calculer un prix de trois façons, fixer le prix de ce qui n'a pas de formule, et savoir à quel point un arbre, une grille ou une simulation se trompe, et pourquoi. Sauf indication contraire dans l'énoncé, le modèle est celui de Black-Scholes, avec un taux constant, une volatilité constante et sans dividendes.
\end{ivtrains}

\section{Échauffements}

\begin{question}[$\star$ \iqroles{QR, QT, QD} \iqfirm{toutes} \iqlevel{junior}]\label{iq:iv2:pricing-and-numerics:1}
Écrivez le prix de Black-Scholes d'un call. Avec un sous-jacent et un strike à 100, une volatilité de 20~\%, un an et un taux nul, donnez le prix de tête par une règle empirique, puis dites quelle est la valeur exacte. Faites de même à 30~\% et pour trois mois.
\end{question}

\begin{question}[$\star$ \iqroles{QR, QT, QD} \iqfirm{toutes} \iqlevel{junior}]\label{iq:iv2:pricing-and-numerics:2}
Deux actions valent toutes deux 100, avec la même volatilité. Tout le monde s'attend à ce que l'une rapporte 20~\% par an, et l'autre 2~\%. Le taux est le même pour les deux. Laquelle a le call à un an le plus cher dans le modèle de Black-Scholes~? Pourquoi un trader sent-il que cela doit être faux, et que répondez-vous~?
\end{question}

\begin{question}[$\star$ \iqroles{QR, QT, QD} \iqfirm{toutes} \iqlevel{junior}]\label{iq:iv2:pricing-and-numerics:3}
Une action à 100 passe demain à 120 ou à 90, et le taux est nul. Pricez un call de strike 100 en construisant un portefeuille d'action et de cash qui rapporte la même chose dans les deux états. Quelle est la probabilité risque-neutre de monter~? Que change-t-il si la probabilité réelle de monter est de 80~\%~?
\end{question}

\begin{question}[$\star$ \iqroles{QR, QD, Risque} \iqfirm{toutes} \iqlevel{junior}]\label{iq:iv2:pricing-and-numerics:4}
Votre payoff a un écart type de 15~\$. Avec 10\,000 trajectoires de Monte Carlo, quelle est l'erreur type du prix~? Combien de trajectoires faut-il pour une erreur type de 1~cent~?
\end{question}

\begin{question}[$\star$ \iqroles{QR, QT, QD} \iqfirm{toutes} \iqlevel{junior}]\label{iq:iv2:pricing-and-numerics:5}
Un call d'un an vaut 8~\$. Le même call, qui disparaît si l'action traite un jour sous 90, vaut 6~\$. Combien vaut le même call s'il n'apparaît qu'une fois l'action passée sous 90~? Pourquoi cela est-il vrai dans n'importe quel modèle~? Vers quoi tend le 6~\$ si la barrière descend vers zéro~?
\end{question}

\begin{question}[$\star$ \iqroles{QR, QT, QD} \iqfirm{toutes} \iqlevel{junior}]\label{iq:iv2:pricing-and-numerics:6}
Un contrat paie 1~\$ si une action finit au-dessus de 100 dans un an, et il vaut 0.46~\$. Le taux est de 5~\%, capitalisé en continu. Combien vaut le contrat qui paie 1~\$ si l'action finit en dessous de 100~? Quelle hypothèse faut-il faire sur l'action~? Quelle volatilité avez-vous utilisée~?
\end{question}

\begin{question}[$\star$ \iqroles{QR, QT, QD} \iqfirm{toutes} \iqlevel{junior}]\label{iq:iv2:pricing-and-numerics:7}
Une option à la monnaie d'un an paie sur la moyenne continue d'une action sur l'année, et non sur son cours final. L'action vaut 100, avec une volatilité de 30~\% et un taux nul. Est-elle plus chère ou moins chère que le call ordinaire~? Quelle volatilité doit-on utiliser à peu près pour la moyenne, et quel est le prix selon la règle empirique~?
\end{question}

\begin{question}[$\star$ \iqroles{QD, QR} \iqfirm{toutes} \iqlevel{junior}]\label{iq:iv2:pricing-and-numerics:8}
Vous construisez un arbre binomial à 20 pas de temps. Combien de nœuds possède-t-il si monter puis descendre aboutit au même prix que descendre puis monter~? Combien s'il n'en est pas ainsi~? Pourquoi cela compte-t-il pour les options sans barrière et indépendantes du chemin, et quelles options perdent cette économie~?
\end{question}

\begin{question}[$\star$ \iqroles{QR, QD} \iqfirm{toutes} \iqlevel{junior}]\label{iq:iv2:pricing-and-numerics:9}
Vous utilisez des trajectoires antithétiques~: chaque tirage normal $Z$ est utilisé avec $-Z$. Que se passe-t-il lorsque vous estimez $E[Z]$~? Et pour $E[Z^2]$~? Pour quels types de payoff l'astuce aide-t-elle~?
\end{question}

\begin{question}[$\star$ \iqroles{QR, QD} \iqfirm{toutes} \iqlevel{junior}]\label{iq:iv2:pricing-and-numerics:10}
Vous disposez de deux normales standard indépendantes, $Z_1$ et $Z_2$. Construisez deux normales de corrélation 0.6. Quelle valeur obtenez-vous pour la seconde si $Z_1 = 1$ et $Z_2 = -1$~? Quel est le rapport avec un facteur de Cholesky~?
\end{question}

\begin{question}[$\star$ \iqroles{QR, QT, QD, Risque} \iqfirm{toutes} \iqlevel{junior}]\label{iq:iv2:pricing-and-numerics:11}
Une action vaut 100, ne verse aucun dividende, et le taux est de 5~\%, capitalisé en continu. Quel est le prix forward juste à un an~? Vous détenez un contrat d'achat à 100 dans un an. Combien vaut-il aujourd'hui~? Pourquoi n'avez-vous pas besoin de volatilité~?
\end{question}

\begin{question}[$\star$ \iqroles{QD, QR} \iqfirm{toutes} \iqlevel{junior}]\label{iq:iv2:pricing-and-numerics:12}
Vous résolvez l'équation de pricing avec une grille de différences finies explicite, en fonction du temps et d'un prix. Vous divisez par deux le pas de prix. De quel facteur le travail augmente-t-il, et pourquoi pas simplement de deux~?
\end{question}

\begin{question}[$\star$ \iqroles{QD, QR} \iqfirm{toutes} \iqlevel{junior}]\label{iq:iv2:pricing-and-numerics:13}
Vous savez tirer des nombres aléatoires uniformes. Comment tirez-vous une exponentielle de paramètre 2~? Quelle valeur obtenez-vous pour $U = 0.5$~? Quelle est la méthode générale, et pourquoi n'est-ce pas la façon habituelle de tirer une normale~?
\end{question}

\section{Calcul de prix}

\begin{question}[$\star\star$ \iqroles{QR, QT, QD} \iqfirm{toutes} \iqlevel{junior}]\label{iq:iv2:pricing-and-numerics:14}
Une action vaut 110, avec une volatilité de 30~\%, six mois avant l'échéance et un taux nul. Sous la probabilité risque-neutre, quelle est la probabilité qu'elle finisse au-dessus de 100~? Établissez-la à partir de la loi du logarithme du cours, et ne vous contentez pas d'écrire $\Phi(d_2)$.
\end{question}

\begin{question}[$\star\star$ \iqroles{QR, QD} \iqfirm{toutes} \iqlevel{junior, confirmé}]\label{iq:iv2:pricing-and-numerics:15}
Dérivez le prix du call de Black-Scholes comme une espérance risque-neutre, avec un taux nul. En particulier, comment traitez-vous $E[S_T \mathbf{1}\{S_T > K\}]$~? Vérifiez-le pour une action de 110~\$, un strike de 100, une volatilité de 30~\% et six mois.
\end{question}

\begin{question}[$\star\star$ \iqroles{QR, QD} \iqfirm{toutes} \iqlevel{junior, confirmé}]\label{iq:iv2:pricing-and-numerics:16}
Dérivez l'équation de Black-Scholes avec un portefeuille couvert en delta. Où est passé le rendement espéré de l'action~? Pour un call d'un an à la monnaie sur une action de 100~\$, avec une volatilité de 20~\% et un taux nul, le gamma vaut 0.0198. Combien doit valoir le theta~?
\end{question}

\begin{question}[$\star\star$ \iqroles{QR, QD} \iqfirm{hedge fund, banque, teneur de marché} \iqlevel{junior, confirmé}]\label{iq:iv2:pricing-and-numerics:17}
Un contrat paie $S_T^2$ dans un an, sur une action à 100~\$ avec une volatilité de 20~\% et un taux nul. Résolvez son prix en cherchant une solution de la forme $V = S^2 g(t)$ dans l'équation de Black-Scholes. Quelle est $g$, et quelle est la valeur~? Vérifiez-le par une espérance.
\end{question}

\begin{question}[$\star\star$ \iqroles{QR, QT, QD} \iqfirm{toutes} \iqlevel{junior}]\label{iq:iv2:pricing-and-numerics:18}
Montrez qu'un call se décompose en une option asset-or-nothing, qui paie l'action si elle finit au-dessus du strike, et une option cash-or-nothing. Que vaut chacune pour une action et un strike de 100, une volatilité de 20~\%, un taux nul et un an~? Pourquoi la valeur de la première n'est-elle pas $S$ multipliée par la probabilité de la seconde~?
\end{question}

\begin{question}[$\star\star$ \iqroles{QR, QT, QD} \iqfirm{banque, hedge fund, teneur de marché} \iqlevel{junior, confirmé}]\label{iq:iv2:pricing-and-numerics:19}
Avec un modèle normal (Bachelier) à la place d'un modèle lognormal, une action à 100 a une volatilité en dollars de 20~\$ par an. Avec un taux nul, combien vaut le call à un an à la monnaie~? Comparez-le avec un modèle lognormal à 20~\%. Quelles volatilités lognormales implicites le modèle normal donne-t-il pour des strikes de 80 et de 120~? Pourquoi présente-t-il un skew~?
\end{question}

\begin{question}[$\star\star$ \iqroles{QR, QT} \iqfirm{toutes} \iqlevel{junior, confirmé}]\label{iq:iv2:pricing-and-numerics:20}
Une action sans drift suit un mouvement brownien, avec une volatilité en dollars de 20~\$ par an, à partir de 100. Quelle est la probabilité qu'elle touche 120 en un an~? Utilisez l'argument de réflexion. Si vous ne regardez que les cours de clôture, sur 252 jours, quelle est à peu près la probabilité~?
\end{question}

\begin{question}[$\star\star$ \iqroles{QR, QT, QD} \iqfirm{banque, hedge fund, teneur de marché} \iqlevel{junior, confirmé}]\label{iq:iv2:pricing-and-numerics:21}
Un contrat paie 1~\$ la première fois qu'une action atteint 110, et il n'expire jamais. L'action vaut 100, c'est une martingale lognormale avec une volatilité de 20~\%, et le taux est nul. Combien vaut-il~? Maintenant il dure un an~: la règle empirique d'un desk consiste à prendre le double du digital européen à 110. Combien cela fait-il, et quelle est la réponse exacte~?
\end{question}

\begin{question}[$\star\star$ \iqroles{QR, QT} \iqfirm{hedge fund, teneur de marché, banque} \iqlevel{junior, confirmé}]\label{iq:iv2:pricing-and-numerics:22}
Un contrat paie $\ln(S_T/S_0)$ à l'échéance. L'action a une volatilité de 20~\%, un taux nul et un an devant elle. Combien vaut-il~? Pourquoi est-il négatif alors que l'action est un pari équitable~? Que obtenez-vous si vous vendez le contrat et détenez $1/S$ actions, en rééquilibrant au fil de l'eau~?
\end{question}

\begin{question}[$\star\star$ \iqroles{QR, QT, QD} \iqfirm{banque, hedge fund, teneur de marché} \iqlevel{junior, confirmé}]\label{iq:iv2:pricing-and-numerics:23}
Dans six mois, on vous donnera un call à la monnaie de six mois, dont le strike est fixé au cours de l'action ce jour-là. L'action vaut 100, avec une volatilité de 20~\% et un taux nul. Combien vaut-il aujourd'hui~? Pourquoi le prix ne dépend-il pas du niveau de l'action dans six mois~?
\end{question}

\begin{question}[$\star\star$ \iqroles{QR, QT, QD} \iqfirm{banque, hedge fund, teneur de marché} \iqlevel{junior, confirmé}]\label{iq:iv2:pricing-and-numerics:24}
Vous détenez un contrat qui vous laisse décider dans six mois si c'est un call ou un put, tous deux de strike 100 et tous deux échéant dans un an. L'action vaut 100, avec une volatilité de 20~\% et un taux nul. Montrez qu'il vaut autant qu'un call d'un an plus un put de six mois. Quelle est la valeur~? Comment cela se compare-t-il à un straddle d'un an~?
\end{question}

\section{Discrétisation}

\begin{question}[$\star\star$ \iqroles{QD, QR} \iqfirm{toutes} \iqlevel{junior, confirmé}]\label{iq:iv2:pricing-and-numerics:25}
Vous simulez une action lognormale par Euler, $S_{i+1} = S_i(1 + \sigma\sqrt{\Delta t}\,Z)$, avec $\sigma = 100~\%$ sur un an et quatre pas, et sans drift. Quelle est la probabilité d'un prix négatif à un pas~? Que vaut $E[S_T^2]/S_0^2$ dans ce schéma, et quelle est la vraie valeur~? Que devriez-vous simuler à la place~?
\end{question}

\begin{question}[$\star\star$ \iqroles{QD, QR} \iqfirm{toutes} \iqlevel{junior, confirmé}]\label{iq:iv2:pricing-and-numerics:26}
Donnez les paramètres de l'arbre binomial standard pour une volatilité de 20~\%, avec un taux nul. Pour $\Delta t = 0.25$, quels sont $u$, $d$ et la probabilité d'un mouvement à la hausse~? Pour un call à un an à la monnaie, la valeur exacte est 7.9656. L'arbre donne 7.7692 avec 10 pas, 8.1472 avec 11, 7.9457 avec 100 et 7.9852 avec 101. Pourquoi oscille-t-il~? Quelle solution peu coûteuse ces chiffres suggèrent-ils~?
\end{question}

\begin{question}[$\star\star$ \iqroles{QD, QR} \iqfirm{toutes} \iqlevel{junior, confirmé}]\label{iq:iv2:pricing-and-numerics:27}
Vous utilisez un schéma explicite aux différences finies sur des prix allant de 0 à 200, par pas de 1~\$, pour une action à 20~\% de volatilité et un taux nul. Quel est le plus grand pas de temps stable, et combien de pas faut-il pour un an~? Que se passe-t-il si vous dépassez un peu ce pas~? Pourquoi cela justifie-t-il de passer à un schéma implicite~?
\end{question}

\begin{question}[$\star\star$ \iqroles{QR, QD, Risque} \iqfirm{toutes} \iqlevel{junior, confirmé}]\label{iq:iv2:pricing-and-numerics:28}
On vous donne une corrélation de 0.9 entre A et B, et de 0.9 entre A et C, et quelqu'un fixe la corrélation de B et C à zéro. Qu'est-ce qui ne va pas lorsque vous prenez une décomposition de Cholesky~? Montrez-le sur un exemple chiffré. Quelle est la corrélation la plus faible entre B et C qui soit admissible~? Comment corrigeriez-vous en pratique une matrice de ce type~?
\end{question}

\begin{question}[$\star\star$ \iqroles{QD, QR, QT} \iqfirm{toutes} \iqlevel{junior, confirmé}]\label{iq:iv2:pricing-and-numerics:29}
Un call d'un an à la monnaie sur une action à 100~\$, avec un taux nul, se négocie à 12~\$. Trouvez la volatilité implicite en partant d'une première estimation issue de la règle empirique, puis par la méthode de Newton. Combien d'itérations vous faut-il pour obtenir six chiffres exacts~? Dans quels cas la méthode de Newton échoue-t-elle en pratique, et quel est le garde-fou~?
\end{question}

\begin{question}[$\star\star$ \iqroles{QD, QR, Risque} \iqfirm{toutes} \iqlevel{junior, confirmé}]\label{iq:iv2:pricing-and-numerics:30}
Vous estimez le delta d'un call à la monnaie sur une action à 100~\$, avec une volatilité de 20~\%, un an et un taux nul, en décalant le cours de l'action de 1~\$ vers le haut et de 1~\$ vers le bas. Vous utilisez 100\,000 trajectoires. L'écart type du payoff est d'environ 13.15~\$. Quelle est l'erreur type du delta si les trajectoires haussières et baissières utilisent des nombres aléatoires différents~? Et s'ils utilisent les mêmes nombres aléatoires~? Quel est, à peu près, l'effet de ce choix~?
\end{question}

\begin{question}[$\star\star$ \iqroles{QR, QT, QD} \iqfirm{toutes} \iqlevel{junior, confirmé}]\label{iq:iv2:pricing-and-numerics:31}
Un panier comporte des poids égaux dans deux actions à 100, toutes deux avec une volatilité de 30~\% et une corrélation de 0.5. Quelle est la volatilité du panier~? Évaluez un call à un an à la monnaie sur ce panier, avec un taux nul, en utilisant une seule loi lognormale. Comment vérifieriez-vous ce résultat~? Quelle est la valeur du call avec une corrélation de 0 et avec une corrélation de 1~?
\end{question}

\begin{question}[$\star\star$ \iqroles{QR, QD} \iqfirm{toutes} \iqlevel{junior, confirmé}]\label{iq:iv2:pricing-and-numerics:32}
Vous évaluez un call à la monnaie par simulation, et vous utilisez le cours final de l'action comme variable de contrôle, puisque vous en connaissez la moyenne. Si la corrélation entre le payoff du call et l'action vaut 0.89, de combien la variance diminue-t-elle~? Quel est le coefficient optimal, et comment l'estimez-vous~? Comment la corrélation évoluerait-elle pour un call profondément dans la monnaie~?
\end{question}

\section{Exotiques et schémas}

\begin{question}[$\star\star\star$ \iqroles{QR, QD} \iqfirm{toutes} \iqlevel{junior, confirmé}]\label{iq:iv2:pricing-and-numerics:33}
Une option paie sur la moyenne géométrique continue d'une action sur un an. L'action et le prix d'exercice valent 100, avec une volatilité de 30~\% et un taux nul. Montrez que la moyenne est lognormale, et déterminez sa volatilité et son forward. Quel est le prix~? Pourquoi existe-t-il une formule fermée ici, et pas pour une moyenne arithmétique~?
\end{question}

\begin{question}[$\star\star\star$ \iqroles{QR, QD} \iqfirm{toutes} \iqlevel{junior, confirmé}]\label{iq:iv2:pricing-and-numerics:34}
Une option paie sur la moyenne arithmétique continue d'une action sur un an. L'action et le prix d'exercice sont à 100, avec une volatilité de 30~\% et un taux nul. Ajustez une loi lognormale aux deux premiers moments. Le moment d'ordre deux de la moyenne continue est $E[A^2] = 2S^2(e^{\sigma^2T} - 1 - \sigma^2T)/(\sigma^2T)^2$. Quel est le prix~? Comment le vérifier par simulation avec peu de trajectoires, et à quoi vous sert une option sur la moyenne géométrique~? De quel facteur l'erreur type diminue-t-elle si les deux payoffs ont une corrélation de 0.9992~?
\end{question}

\begin{question}[$\star\star\star$ \iqroles{QR, QT, QD} \iqfirm{toutes} \iqlevel{junior, confirmé}]\label{iq:iv2:pricing-and-numerics:35}
Une option vous permet d'échanger l'action B contre l'action A à l'échéance, avec un payoff de $(S_A - S_B)^+$. Les deux actions valent 100. Leurs volatilités sont de 30~\% et 20~\%, la corrélation vaut 0.5, et il s'agit d'une option à un an. Quelle est sa valeur~? Montrez pourquoi le taux n'intervient pas. Quelle est l'astuce~? Vérifiez-la dans la limite où B n'a aucune volatilité.
\end{question}

\begin{question}[$\star\star\star$ \iqroles{QR, QT} \iqfirm{banque, hedge fund, teneur de marché} \iqlevel{junior, confirmé}]\label{iq:iv2:pricing-and-numerics:36}
Un call à un an de strike 100 s'éteint si l'action se négocie à 100 ou moins à un instant quelconque. L'action vaut 110, avec un taux nul. Quelle est sa valeur, dans un modèle avec une martingale continue quelconque pour l'action~? Pourquoi la volatilité n'intervient-elle pas~? Comparez-la à un call ordinaire à 30~\% de volatilité, et dites ce que vaut la différence pour son détenteur.
\end{question}

\begin{question}[$\star\star\star$ \iqroles{QR, QT, QD} \iqfirm{banque, hedge fund, teneur de marché} \iqlevel{junior, confirmé}]\label{iq:iv2:pricing-and-numerics:37}
Un call à un an de strike 100 s'éteint si l'action atteint 130 à un moment quelconque. L'action vaut 100, avec un taux nul. Valorisez-le par la méthode des images, en prenant 20~\%, 15~\% et 25~\% comme volatilités. Quel est le signe de son vega~? Comment cela se compare-t-il à un call ordinaire, et pourquoi~?
\end{question}

\begin{question}[$\star\star\star$ \iqroles{QR, QD, QT} \iqfirm{banque, hedge fund} \iqlevel{junior, confirmé}]\label{iq:iv2:pricing-and-numerics:38}
Un call up-and-out à un an sur une action à 100 a un strike de 100, une barrière à 130, une volatilité de 20~\% et un taux nul. La barrière n'est observée qu'à la clôture, sur 252 jours. Son prix continu vaut 2.97. Le prix quotidien devrait-il être plus élevé ou plus bas~? Quel est le décalage usuel de la barrière qui corrige la formule, et que donne-t-il~? Quelle est l'ampleur de l'effet sur le niveau de la barrière~?
\end{question}

\begin{question}[$\star\star\star$ \iqroles{QR, QD} \iqfirm{banque, hedge fund} \iqlevel{junior, confirmé}]\label{iq:iv2:pricing-and-numerics:39}
Vous valorisez un call up-and-out à un an, à barrière observée en continu à 130, avec l'action et le prix d'exercice à 100, une volatilité de 20~\% et un taux nul. Avec des pas mensuels de Monte-Carlo, vous obtenez 3.73. La bonne réponse est 2.97. Quel est le biais~? Comment le corrigeriez-vous sans ajouter de pas~? Étant donné un mois qui commence à 120 et se termine à 125, avec une volatilité de 20~\%, quelle est la probabilité qu'il ait touché 130 entre-temps~?
\end{question}

\begin{question}[$\star\star\star$ \iqroles{QR, QD, Risque} \iqfirm{toutes} \iqlevel{confirmé}]\label{iq:iv2:pricing-and-numerics:40}
Vous voulez le delta d'un call par simulation, pour une action à 100~\$ et un prix d'exercice de 100~\$, avec une volatilité de 20~\%, un an et un taux nul. Dérivez l'estimateur pathwise et l'estimateur par rapport de vraisemblance. Quelles sont leurs moyennes et leurs écarts types par trajectoire~? Lequel utiliseriez-vous pour un digital, et pourquoi~?
\end{question}

\begin{question}[$\star\star\star$ \iqroles{QR, QT} \iqfirm{toutes} \iqlevel{junior, confirmé}]\label{iq:iv2:pricing-and-numerics:41}
Une action suit un mouvement brownien sans drift et une volatilité en dollars de 20~\$ par an. Un contrat verse la différence entre son cours le plus haut de l'année et son cours initial, et il est observé en continu. Quelle est sa valeur, par le principe de réflexion~? Comment cela se compare-t-il à un straddle à la monnaie dans le même modèle~? Qu'attendriez-vous si l'on ne vérifiait qu'une fois par jour~?
\end{question}

\begin{question}[$\star\star\star$ \iqroles{QR, QT, QD} \iqfirm{teneur de marché, banque, hedge fund} \iqlevel{junior, confirmé}]\label{iq:iv2:pricing-and-numerics:42}
Un digital de 1~\$ paie au-dessus de 100 dans un an. L'action vaut 100 et le taux est nul. La volatilité implicite à la monnaie est de 20~\%, le vega vaut 39.7 par unité de volatilité, et le smile diminue de 0.2 point de volatilité pour chaque 1~\$ de prix d'exercice. Quelle est la valeur du digital sans skew, et avec le skew~? Dérivez la correction. Quel est le signe pour la plupart des indices boursiers~? De combien cela modifie-t-il la probabilité~?
\end{question}

\begin{question}[$\star\star\star$ \iqroles{QR, QD} \iqfirm{toutes} \iqlevel{confirmé}]\label{iq:iv2:pricing-and-numerics:43}
Vous voulez la probabilité qu'une loi normale centrée réduite dépasse 4. Avec un million de tirages Monte-Carlo classiques, quelle est l'erreur type relative~? Décalez les tirages de 4 et repondérez-les. Quel est le poids, et quelle est désormais l'erreur type relative~? Que peut-il mal se passer avec un décalage lorsqu'on le fait en grande dimension~?
\end{question}

\begin{question}[$\star\star\star$ \iqroles{QR, QD} \iqfirm{toutes} \iqlevel{confirmé}]\label{iq:iv2:pricing-and-numerics:44}
Vous valorisez un call à un an à la monnaie avec un tirage normal par trajectoire, soit $2^{14}$ trajectoires, et une volatilité de 20~\%. Comparez l'erreur du Monte-Carlo classique à celle d'une suite de Sobol randomisée, mesurée sur vingt graines. Pourquoi les points de Sobol font-ils mieux~? Dans quelles situations perdez-vous l'essentiel de cet avantage~?
\end{question}

\begin{question}[$\star\star\star$ \iqroles{QR, QD} \iqfirm{toutes} \iqlevel{junior, confirmé}]\label{iq:iv2:pricing-and-numerics:45}
Vous valorisez un put américain à un an sur une action, de prix d'exercice 100, avec une volatilité de 20~\% et un taux de 5~\%. Un arbre très fin donne 6.0903, et le put européen vaut 5.5735. Un arbre à 100 pas donne 6.0824, un arbre à 101 pas 6.1046 et un arbre à 200 pas 6.0864. L'arbre européen à 100 pas donne 5.5536. Comparez la valeur brute, la moyenne des arbres à 100 et 101 pas, la variable de contrôle européenne et le doublement du nombre de pas avec extrapolation. Quelle est la meilleure approche ici~? Quel est le choix usuel dans les manuels~? Quel est le cours de l'action en dessous duquel vous exercez aujourd'hui~?
\end{question}

\begin{question}[$\star\star\star$ \iqroles{QD, QR} \iqfirm{toutes} \iqlevel{confirmé}]\label{iq:iv2:pricing-and-numerics:46}
Vous résolvez un call à cinq jours de l'échéance, $T = 0.02$, par Crank-Nicolson sur une grille de prix fine, avec un pas de 0.25~\$ et cinq pas de temps. Le prix est assez proche, mais le gamma au prix d'exercice ressort à $-1.36$. Le vrai gamma vaut 0.141. Que se passe-t-il, alors même que le schéma est inconditionnellement stable~? Quelle est la parade usuelle~? Quel est son effet sur ce gamma~?
\end{question}

\begin{question}[$\star\star\star$ \iqroles{QR, QD, Risque} \iqfirm{toutes} \iqlevel{junior, confirmé}]\label{iq:iv2:pricing-and-numerics:47}
Déterminez l'écart type exact du payoff d'un call d'un an sur une action à 100~\$, avec une volatilité de 20~\% et un taux nul, en calculant $E[((S_T - K)^+)^2]$ sous forme fermée, pour des prix d'exercice de 100 et de 130. Combien de trajectoires vous faut-il pour une erreur type d'un centime~? Comment l'erreur se compare-t-elle au prix pour chaque prix d'exercice~? Qu'est-ce que cela dit de la valorisation des options de queue~?
\end{question}

\section{Trois à construire}

\begin{question}[$\star\star\star\star$ \iqroles{QR, QD} \iqfirm{banque, hedge fund, teneur de marché} \iqlevel{confirmé}]\label{iq:iv2:pricing-and-numerics:48}
Décrivez, étape par étape, la méthode de régression de Longstaff-Schwartz pour un put bermudéen à 12 dates d'exercice mensuelles, de prix d'exercice 100 sur une action à 100~\$, avec une volatilité de 20~\% et un taux de 5~\%. Un arbre donne 6.04. Que donne la méthode avec 100\,000 trajectoires, et avec 400~? Dans quel sens est-elle biaisée lorsque vous réutilisez les trajectoires d'ajustement, et dans quel sens sur des trajectoires fraîches~? Pourquoi ne régressez-vous que sur les trajectoires dans la monnaie~?
\end{question}

\begin{question}[$\star\star\star\star$ \iqroles{QD, QR} \iqfirm{banque, hedge fund, teneur de marché} \iqlevel{confirmé}]\label{iq:iv2:pricing-and-numerics:49}
Formulez un put américain comme problème aux différences finies sur une grille de prix de l'action, pour une action à 100~\$ de prix d'exercice 100, avec une volatilité de 20~\%, un taux de 5~\% et un an avant l'échéance. Écrivez la condition qui vaut en chaque point, et la manière de la résoudre à chaque pas de temps. Quel résultat obtenez-vous avec une grille de 1~\$ et 100 pas de temps, à comparer à 6.0903~? Où se situe aujourd'hui la frontière d'exercice~? Que se passe-t-il si vous résolvez le système linéaire puis le tronquez au niveau du payoff~?
\end{question}

\begin{question}[$\star\star\star\star$ \iqroles{QR, QD} \iqfirm{hedge fund, banque, teneur de marché} \iqlevel{confirmé}]\label{iq:iv2:pricing-and-numerics:50}
Étant donné la fonction caractéristique du logarithme du prix dans un modèle, $\phi(u) = E[e^{iu\ln S_T}]$, dérivez le prix du call à partir de deux probabilités, et écrivez-les sous forme d'intégrales. Testez-le avec Black-Scholes pour des prix d'exercice de 100 et 120, sur une action à 100~\$, avec une volatilité de 20~\%, un taux nul et un an. Que donne-t-il, et combien perdez-vous à cause de la borne supérieure d'intégration~? Pourquoi le faire si une formule fermée existe~?
\end{question}
